% =============================================================================
% Appendix-style section: Confounder Validation of the Insufficiency Probe
% Status: draft, section-only (no \documentclass / preamble)
% Last updated: 2026-05-13
%
% Required packages in the main paper preamble:
%   \usepackage{booktabs}
%   \usepackage{array}
%   \usepackage{amsmath}
% =============================================================================

\section{Confounder Validation of the Insufficiency Probe}
\label{sec:confounder}

\subsection{Motivation}

We train a linear probe on residual-stream activations at $t{=}0$
(the end-of-prompt token, before any generation) to classify whether a
math word problem is logically insufficient. The probe achieves high
held-out F1 across $21$ models on UMWP (mean $0.87$) and TreeCut (mean
$0.79$). A natural concern is whether this F1 reflects a genuine
semantic signal or merely a surface artifact such as token length,
numeric token count, problem complexity, or insufficiency-indicating
lexical cues. We address this through three complementary defenses,
applied at $t{=}0$ on the solving prompt that matches the probe's
training distribution. For clarity of exposition, we report results on
six representative models spanning $1.5$B--$70$B parameters and four
families (Qwen, LLaMA, Gemma, DeepSeek-R1). Full $21$-model results are
deferred to a separate table.

\subsection{Defense 1: Minimal Contrastive Pairs}
\label{sec:d1}

\paragraph{Construction.}
For each dataset we sample $250$ solvable questions $q^+$ from the test
set and use a powerful LLM to produce a minimally edited unsolvable
counterpart $q^-$ by substituting a single numerical value with a vague
placeholder (e.g.\ ``$27$ files'' $\to$ ``some files''). A second LLM
judge verifies (i) the edit is minimal ($\le 3$ words changed) and
(ii) $q^-$ is genuinely unsolvable under reasonable assumptions.
Verified pair counts and mean token edit distances:
$239/250$ verified, edit distance $1.06$ on UMWP;
$232/250$ verified, edit distance $1.00$ on TreeCut.

\paragraph{What it controls.}
$q^+$ and $q^-$ share length, scenario, and almost all lexical content.
A probe that reduces to length, scenario, or word distribution cannot
systematically assign higher $P(\text{insuff})$ to $q^-$.

\paragraph{Reported metric.}
We report
\[
\mathrm{Mean\_Gap} \;=\; \mathbb{E}_{(q^+,q^-)}\!\left[\,
   P(\text{insuff}\mid q^-) - P(\text{insuff}\mid q^+)\,\right],
\]
together with the raw class-conditional probabilities
$\overline{P(\text{insuff}\mid q^+)}$ and
$\overline{P(\text{insuff}\mid q^-)}$ to expose probe degeneracy
(saturation at $0$ or $1$).

\begin{table}[h]
\centering
\small
\caption{D1 results. On UMWP, all six probes assign substantially higher
insufficiency probability to $q^-$ than $q^+$
($\mathrm{Mean\_Gap} \in [0.52, 0.86]$). On TreeCut, only the Gemma
models retain a non-trivial gap; the remaining four either saturate
near $0$ ($\overline{P(\cdot)} < 0.01$) or near $1$
($\overline{P(\cdot)} \approx 0.88$).}
\label{tab:d1}
\begin{tabular}{llrrrr}
\toprule
Dataset & Model & Layer
        & $\overline{P}\!\mid\!q^+$
        & $\overline{P}\!\mid\!q^-$
        & $\mathrm{Mean\_Gap}$ \\
\midrule
\multirow{6}{*}{UMWP}
 & Qwen2.5-Math-1.5B-Instruct      & 18 & 0.24 & 0.94 & \textbf{0.70} \\
 & Qwen2.5-Math-7B                 & 20 & 0.02 & 0.77 & \textbf{0.75} \\
 & Meta-Llama-3.1-8B               & 19 & 0.31 & 0.83 & \textbf{0.52} \\
 & gemma-3-12b-it                  & 26 & 0.18 & 0.94 & \textbf{0.76} \\
 & gemma-3-27b-it                  & 34 & 0.06 & 0.92 & \textbf{0.86} \\
 & DeepSeek-R1-Distill-Llama-70B   & 31 & 0.26 & 0.96 & \textbf{0.70} \\
\midrule
\multirow{6}{*}{TreeCut}
 & Qwen2.5-Math-1.5B-Instruct      & 16 & 0.88 & 0.93 & 0.06 \\
 & Qwen2.5-Math-7B                 & 15 & 0.88 & 0.89 & 0.01 \\
 & Meta-Llama-3.1-8B               &  9 & 0.00 & 0.00 & 0.00 \\
 & gemma-3-12b-it                  & 25 & 0.14 & 0.54 & \textbf{0.41} \\
 & gemma-3-27b-it                  & 28 & 0.21 & 0.67 & \textbf{0.46} \\
 & DeepSeek-R1-Distill-Llama-70B   & 33 & 0.00 & 0.03 & 0.03 \\
\bottomrule
\end{tabular}
\end{table}

\paragraph{Interpretation.}
On UMWP, a single-word edit that preserves length and lexical content
flips the probe by a mean of $0.70$ probability mass. The probe is
responding to semantic content, not surface features. On TreeCut the
defense holds only for the two Gemma models; for the other four,
probe saturation (e.g.\ $\overline{P} \approx 0.88$ for both classes on
Qwen2.5-Math-7B) makes the test uninformative for those models. We do
not claim a clean semantic signal for TreeCut on saturated probes.

\subsection{Defense 2: Adversarial Lexical Injection}
\label{sec:d2}

\paragraph{Construction.}
We mine the top TF-IDF cue tokens from the insufficient class and
inject them into solvable test questions. A probe relying on these
keywords should produce a false-positive flip; a probe reading
semantics should remain stable.

\paragraph{Filter.}
We evaluate only on (model, dataset) pairs where the baseline
false-positive rate is below $0.5$. Above this threshold the probe is
already saturated for the dataset and the FP rate cannot meaningfully
grow, making the test degenerate.

\paragraph{Reported metric.}
$\mathrm{FP\_Excess} = \overline{\mathrm{FP\_rate}_{\text{injected}}}
- \mathrm{FP\_rate}_{\text{baseline}}$.
A semantic probe should satisfy $\mathrm{FP\_Excess} \approx 0$.

\begin{table}[h]
\centering
\small
\caption{D2 results, restricted to (model, dataset) pairs with
non-saturated baselines. Cue sets are dataset-specific:
\textit{UMWP}: \{apples, earned, many, girls, games\};
\textit{TreeCut}: \{many, texas, texas bbq, plate, unspecified\}.}
\label{tab:d2}
\begin{tabular}{llrrr}
\toprule
Dataset & Model & Baseline\_FP & Injection\_FP & FP\_Excess \\
\midrule
\multirow{6}{*}{UMWP}
 & Qwen2.5-Math-1.5B-Instruct       & 0.14 & 0.26 & 0.11 \\
 & Qwen2.5-Math-7B                  & 0.02 & 0.06 & 0.04 \\
 & Meta-Llama-3.1-8B                & 0.33 & 0.45 & 0.13 \\
 & gemma-3-12b-it                   & 0.19 & 0.37 & 0.18 \\
 & gemma-3-27b-it                   & 0.06 & 0.11 & 0.05 \\
 & DeepSeek-R1-Distill-Llama-70B    & 0.26 & 0.66 & \textbf{0.39} \\
\midrule
\multirow{4}{*}{TreeCut}
 & Meta-Llama-3.1-8B                & 0.00 & 0.00 & 0.00 \\
 & gemma-3-12b-it                   & 0.08 & 0.11 & 0.03 \\
 & gemma-3-27b-it                   & 0.20 & 0.20 & $-0.01$ \\
 & DeepSeek-R1-Distill-Llama-70B    & 0.00 & 0.01 & 0.00 \\
\bottomrule
\end{tabular}
\end{table}

\paragraph{Interpretation.}
Five of the six UMWP probes show $\mathrm{FP\_Excess} \le 0.18$
(robust to lexical cues). The DeepSeek-R1-Distill-Llama-70B probe shows
$\mathrm{FP\_Excess} = 0.39$, indicating partial keyword sensitivity;
we report this honestly rather than excluding the model. On TreeCut,
the four non-saturated probes all show $\mathrm{FP\_Excess} < 0.05$.
We treat D2 as supplementary corroboration of D1; it cannot bound the
marginal lexical contribution but rules out a pure-keyword probe on
most models.

\subsection{Defense 3: Geometric Separability from Surface Confounders}
\label{sec:d3}

\paragraph{Construction.}
For each model and dataset, we estimate four ``confounder directions''
in residual space by training auxiliary linear regressors mapping
hidden states to (i) token length, (ii) numeric token count,
(iii) Flesch reading-ease proxy, and (iv) presence of insufficient
lexical cues. The probe direction $v_{\text{probe}}$ is the
StandardScaler-corrected logistic-regression weight from the unified
exp10 probe at the same layer.

\paragraph{Reported metric.}
$\cos(v_{\text{probe}}, v_{\text{conf}_i})$ for each confounder $i$.
A value near $0$ indicates the probe direction is orthogonal to that
confounder direction in activation space: the probe is not lying within
the subspace spanned by surface features.

\begin{table}[h]
\centering
\small
\caption{D3 cosines. All $|\cos|$ are below $0.08$. The insufficiency
direction is geometrically near-orthogonal to all four measured surface
directions on both datasets. (TreeCut does not include a lexical
confounder direction because TF-IDF cue mining was too noisy to
produce a single direction.)}
\label{tab:d3}
\begin{tabular}{llrrrr}
\toprule
Dataset & Model
   & $\cos$\,length
   & $\cos$\,numcount
   & $\cos$\,complexity
   & $\cos$\,lexical \\
\midrule
\multirow{6}{*}{UMWP}
 & Qwen2.5-Math-1.5B-Instruct      & $\phantom{-}0.012$ & $-0.075$ & $\phantom{-}0.017$ & $\phantom{-}0.063$ \\
 & Qwen2.5-Math-7B                 & $-0.027$ & $-0.027$ & $-0.011$ & $\phantom{-}0.030$ \\
 & Meta-Llama-3.1-8B               & $-0.004$ & $-0.017$ & $\phantom{-}0.005$ & $\phantom{-}0.025$ \\
 & gemma-3-12b-it                  & $\phantom{-}0.007$ & $\phantom{-}0.018$ & $\phantom{-}0.051$ & $\phantom{-}0.057$ \\
 & gemma-3-27b-it                  & $\phantom{-}0.065$ & $\phantom{-}0.035$ & $\phantom{-}0.015$ & $\phantom{-}0.011$ \\
 & DeepSeek-R1-Distill-Llama-70B   & $-0.007$ & $\phantom{-}0.007$ & $\phantom{-}0.008$ & $\phantom{-}0.015$ \\
\midrule
\multirow{6}{*}{TreeCut}
 & Qwen2.5-Math-1.5B-Instruct      & $\phantom{-}0.017$ & $\phantom{-}0.013$ & $-0.068$ & --- \\
 & Qwen2.5-Math-7B                 & $\phantom{-}0.002$ & $\phantom{-}0.023$ & $-0.055$ & --- \\
 & Meta-Llama-3.1-8B               & $\phantom{-}0.026$ & $-0.009$ & $-0.024$ & --- \\
 & gemma-3-12b-it                  & $-0.039$ & $\phantom{-}0.021$ & $-0.020$ & --- \\
 & gemma-3-27b-it                  & $\phantom{-}0.020$ & $-0.010$ & $-0.017$ & --- \\
 & DeepSeek-R1-Distill-Llama-70B   & $-0.009$ & $\phantom{-}0.033$ & $-0.006$ & --- \\
\bottomrule
\end{tabular}
\end{table}

\subsection{Confound vs. Natural Correlate}
\label{sec:reframe}

The four features tested by D3 are natural correlates of insufficiency
in math word problems: an insufficient problem typically has one fewer
numerical value (lower numcount), may contain hedging vocabulary
(lexical cues), and may be phrased differently (length, complexity).
These correlations are properties of the input distribution, not
artifacts to control away. Following \citet{arditi2024refusal} on
refusal directions, we adopt the position that the appropriate test
is whether the probe direction \emph{reduces to} surface features,
not whether it correlates with them. D1 (single-word semantic edits
flip the UMWP probe by $\overline{\mathrm{Mean\_Gap}} = 0.71$) and D3
(all $|\cos| < 0.08$ across both datasets) jointly establish that
the probe direction captures a semantic signal that is correlated
with surface features by data construction but is not contained in
their span.

\subsection{What We Claim}

\begin{itemize}
\item On UMWP, the unified probe direction $(i)$ responds to
single-word semantic perturbations
($\mathrm{Mean\_Gap} \in [0.52, 0.86]$ across all six representative
models), $(ii)$ is largely insensitive to lexical-cue injection
($\mathrm{FP\_Excess} \le 0.18$ for five of six models), and
$(iii)$ is geometrically near-orthogonal to length, numeric-count,
complexity, and lexical directions ($|\cos| \le 0.08$). The probe
therefore reads a semantic signal not reducible to the surface
features tested.

\item On TreeCut, D3 (geometric separability) and D2
(lexical robustness on non-saturated probes) hold uniformly. D1
holds only for the two Gemma models in our representative set. The
remaining models exhibit probe saturation that renders D1
uninformative rather than refuted.
\end{itemize}

\subsection{Limitations}

\begin{itemize}
\item TreeCut probing is weaker than UMWP probing by every measure.
Several model probes saturate at $P(\text{insuff})\!\approx\!0$ or
$\approx\!1$ regardless of input, which prevents D1 from
discriminating $q^+$ from $q^-$ on those models. We do not claim a
clean semantic signal for TreeCut on saturated probes.

\item The four confounder directions in D3 (length, numcount,
complexity, lexical) are not exhaustive. Higher-order surface
features such as syntactic structure, discourse markers, or
operator distribution are not tested here.

\item One model (DeepSeek-R1-Distill-Llama-70B) shows partial
lexical-cue sensitivity on UMWP D2
($\mathrm{FP\_Excess} = 0.39$). The semantic and geometric evidence
(D1 $\mathrm{Mean\_Gap} = 0.70$, all D3 cosines $\le 0.015$) supports
a non-pure-keyword probe even for this model, but the lexical
contribution cannot be quantitatively bounded by these defenses.

\item Defenses are applied at $t{=}0$. The robustness of the probe
during chain-of-thought generation is established separately by
matching of Unified and Separate (oracle, retrained per cutoff)
probe F1 across cutoffs ($\mathrm{mean}|\Delta\mathrm{F1}|<0.01$),
which is the subject of the trajectory analysis section.
\end{itemize}

% =============================================================================
% Section: Probe trajectory vs.\ verbalization trajectory
% =============================================================================

\section{Two Findings on the Effect of Chain-of-Thought}
\label{sec:cot-findings}

We compare two trajectories across chain-of-thought positions
$\tau \in \{0, 20, 40, 60, 100\}$\% for every (model, dataset, sample):

\begin{itemize}
\item \textbf{Internal trajectory} (probe F1): the unified linear probe
(Section~\ref{sec:confounder}) evaluated on residual-stream activations
extracted at position $\tau$. This measures the discriminative power of
the internal representation at $\tau$.

\item \textbf{Behavioral trajectory} (forced-decoded verbalization):
truncate the model's CoT at $\tau$, append
\verb|\n\n**Final Answer**\n\boxed{|, and greedy-decode the answer.
A judge grades the result. We report
$\mathrm{Q2\_acc} = P(\text{model writes ``Insufficient''}\,\mid\,\text{insufficient input})$
and
$\mathrm{Q3\_acc} = P(\text{model gives correct number}\,\mid\,\text{sufficient input})$.
\end{itemize}

Both trajectories are computed on the same $21$ models, two datasets
(UMWP, TreeCut), and the same five cutoff positions.
Population means are shown in Table~\ref{tab:traj-popmeans}.

\begin{table}[h]
\centering
\small
\caption{Population-mean trajectories across $21$ models. ``$\Delta$ from 0\%''
rows isolate the direction and magnitude of change at each cutoff.}
\label{tab:traj-popmeans}
\begin{tabular}{lrrrrr}
\toprule
                                 &    0\% &   20\% &   40\% &   60\% &  100\% \\
\midrule
\multicolumn{6}{l}{\emph{UMWP}} \\
\addlinespace[2pt]
Probe F1 (internal)              & 0.871 & 0.744 & 0.768 & 0.791 & 0.869 \\
\quad $\Delta$ from 0\%          & $\phantom{-}0.000$ & $\mathbf{-0.127}$ & $-0.103$ & $-0.080$ & $-0.002$ \\
Q2 abstention (verbalization)    & 0.444 & 0.493 & 0.491 & 0.480 & 0.556 \\
\quad $\Delta$ from 0\%          & $\phantom{-}0.000$ & $\mathbf{+0.049}$ & $+0.048$ & $+0.037$ & $\mathbf{+0.112}$ \\
Q3 solving (verbalization)       & 0.378 & 0.384 & 0.497 & 0.639 & 0.873 \\
\quad $\Delta$ from 0\%          & $\phantom{-}0.000$ & $+0.006$ & $+0.119$ & $+0.261$ & $\mathbf{+0.495}$ \\
\addlinespace[4pt]
\multicolumn{6}{l}{\emph{TreeCut}} \\
\addlinespace[2pt]
Probe F1 (internal)              & 0.785 & 0.618 & 0.618 & 0.652 & 0.722 \\
\quad $\Delta$ from 0\%          & $\phantom{-}0.000$ & $\mathbf{-0.167}$ & $-0.167$ & $-0.133$ & $-0.063$ \\
Q2 abstention (verbalization)    & 0.167 & 0.206 & 0.258 & 0.251 & 0.262 \\
\quad $\Delta$ from 0\%          & $\phantom{-}0.000$ & $\mathbf{+0.040}$ & $+0.091$ & $+0.084$ & $\mathbf{+0.096}$ \\
Q3 solving (verbalization)       & 0.108 & 0.112 & 0.148 & 0.208 & 0.485 \\
\quad $\Delta$ from 0\%          & $\phantom{-}0.000$ & $+0.004$ & $+0.040$ & $+0.100$ & $\mathbf{+0.378}$ \\
\bottomrule
\end{tabular}
\end{table}

\subsection{Finding 1: Probe and verbalization move in opposite directions
at the onset of reasoning}
\label{sec:opposite-direction}

At the $0\!\to\!20\%$ cutoff transition --- the moment when the model
transitions from no CoT to its first chunk of reasoning --- the internal
probe and the Q2 abstention verbalization move in \emph{opposite directions
on the same models with the same data}:

\begin{center}
\small
\begin{tabular}{lrr}
\toprule
                  & UMWP & TreeCut \\
\midrule
$\Delta$ Probe F1 at $0\!\to\!20\%$               & $-0.127$ & $-0.167$ \\
$\Delta$ Q2 verbalization at $0\!\to\!20\%$       & $+0.049$ & $+0.040$ \\
\bottomrule
\end{tabular}
\end{center}

The internal representation loses $13$--$17$ points of insufficiency
discriminability at the same cutoff position where behavioral abstention
gains $4$--$5$ points. The sign reversal is observed on the same
$(\text{model}, \text{dataset}, \text{sample})$ triples; it is not a
between-models artifact.

This is a direct contradiction of the simplest hypothesis that probe and
behavior are aligned readouts of the same internal state. They are not
aligned at the most informative position: the first reasoning step.

\subsection{Finding 2: Chain-of-thought helps solving far more than it
helps abstaining}
\label{sec:asymmetric-cot-gain}

Over the full CoT ($0\!\to\!100\%$), reasoning length improves the model's
ability to produce correct numerical answers (Q3) much more than its
ability to correctly abstain from unsolvable problems (Q2):

\begin{center}
\small
\begin{tabular}{lrrr}
\toprule
Dataset & $\Delta$ Q3 (solving) & $\Delta$ Q2 (abstention) & Ratio \\
\midrule
UMWP    & $+0.495$ & $+0.112$ & $4.4\times$ \\
TreeCut & $+0.378$ & $+0.096$ & $3.9\times$ \\
\bottomrule
\end{tabular}
\end{center}

CoT-induced improvement is $\approx 4\times$ larger for solving than for
abstention. The asymmetry is consistent across both datasets and across
the $21$-model population mean.

This is empirically distinct from the standard claim that ``CoT helps
performance.'' What CoT chiefly helps is the production of answers on
solvable inputs. The capacity to correctly decline an unsolvable input
grows much more slowly with reasoning length. Models do not become
proportionally more epistemically cautious as they think longer; they
become substantially more productive of confident answers.

\subsection{What we claim}

\begin{enumerate}
\item \emph{Internal and behavioral signals diverge at the first reasoning
step.} At cutoff $0\!\to\!20\%$, mean probe F1 falls by $0.127$ (UMWP) /
$0.167$ (TreeCut); mean Q2 verbalization rises by $0.049$ / $0.040$.
Opposite directions on identical $(\text{model}, \text{dataset})$ pairs,
$21$ models, two datasets.

\item \emph{Chain-of-thought has an asymmetric epistemic effect.}
Over the full CoT, solving accuracy gains $+0.495$ (UMWP) / $+0.378$
(TreeCut) while abstention accuracy gains only $+0.112$ / $+0.096$.
Reasoning is approximately four times more beneficial for producing
answers than for declining to answer.
\end{enumerate}

Both claims are population-level statistical observations. They are not
causal at the sample level. We make no claim that the probe trajectory
predicts the verbalization outcome of any individual sample, nor that
manipulating the probe direction will alter verbalization.

\subsection{Limitations}

\begin{itemize}
\item Verbalization is measured by forced-decoded answers after a fixed
suffix. Free-form continuation of the CoT may produce different
abstention rates, especially in models trained to use specific stop
tokens.

\item Q2 sample counts vary across models; on TreeCut several
instruction-tuned models rarely abstain at all
($N_{\text{Q2}}<30$). Population means weight models equally regardless
of effective sample size. Per-model trajectories with very low $N_{\text{Q2}}$
should be read as illustrative.

\item The probe's internal trajectory is itself validated against per-cutoff
oracle probes ($\mathrm{mean}|\Delta\mathrm{F1}|<0.01$ between unified and
separately retrained probes), ruling out a single-probe overfitting
artifact for the cliff-and-recover shape.

\item The opposite-direction observation is robust at the population level
but does not establish that probe F1 \emph{causes}, or is mechanistically
prior to, verbalization. The two trajectories are computed from
the same hidden state stream and are not independent random variables;
their correlation structure at the sample level is the subject of
further analysis.
\end{itemize}

% =============================================================================
% Bibliography entry needed in main paper:
% @article{arditi2024refusal,
%   title={Refusal in language models is mediated by a single direction},
%   author={Arditi, Andy and Obeso, Oscar and Syed, Aaquib and
%           Paleka, Daniel and Panickssery, Nina and
%           Gurnee, Wes and Nanda, Neel},
%   journal={arXiv preprint arXiv:2406.11717},
%   year={2024}
% }
% =============================================================================
